\section{Affine transfer and partition differences}\label{sec:residue-affine}

\subsection{Fourier formulas for the fibers}

\begin{theorem}[Residue fibers]
\label{thm:residue-fibers}
For every \(x\in X_m\), with \(r_x:=V_m(x)\),
\[
  d_m(x)
  =
  \card{\Bigl\{\omega\in\Omega_m:
  \sum_{j=1}^m\omega_jF_{j+1}\equiv r_x \pmod{F_{m+2}}\Bigr\}}.
\]
\end{theorem}

\begin{proof}
This is exactly \Cref{cor:residue-characterization}.
\end{proof}

\begin{theorem}[Fourier inversion]
\label{thm:fourier-inversion}
Let \(\zeta=e^{2\pi i/F_{m+2}}\).  Then for every \(x\in X_m\),
\[
  d_m(x)=\frac{1}{F_{m+2}}
  \sum_{a=0}^{F_{m+2}-1}
  \zeta^{-ar_x}\prod_{j=1}^m(1+\zeta^{aF_{j+1}}).
\]
Consequently, if
\[
  \Sigma_m:=\sum_{x\in X_m}d_m(x)^2,
\]
then
\[
  \Sigma_m=
  \frac{1}{F_{m+2}}
  \sum_{a=0}^{F_{m+2}-1}
  \left|
    \prod_{j=1}^m(1+\zeta^{aF_{j+1}})
  \right|^2.
\]
\end{theorem}

\begin{proof}
In the finite cyclic group \(\ZZ/F_{m+2}\ZZ\),
\[
  \mathbf{1}_{n\equiv r_x}
  =
  \frac{1}{F_{m+2}}
  \sum_{a=0}^{F_{m+2}-1}\zeta^{a(n-r_x)}.
\]
Therefore
\[
  d_m(x)
  =
  \sum_{\omega\in\Omega_m}\mathbf{1}_{N_m(\omega)\equiv r_x}
  =
  \frac{1}{F_{m+2}}
  \sum_{a=0}^{F_{m+2}-1}
  \zeta^{-ar_x}
  \sum_{\omega\in\Omega_m}\zeta^{aN_m(\omega)}.
\]
Since \(N_m(\omega)=\sum_{j=1}^m \omega_jF_{j+1}\),
\[
  \sum_{\omega\in\Omega_m}\zeta^{aN_m(\omega)}
  =
  \prod_{j=1}^m\sum_{\omega_j\in\{0,1\}}\zeta^{a\omega_jF_{j+1}}
  =
  \prod_{j=1}^m(1+\zeta^{aF_{j+1}}),
\]
which proves the first formula.

For the second formula, use that \(d_m(x)\in\RR\), so
\(\Sigma_m=\sum_{x\in X_m}|d_m(x)|^2\).  Multiply the first identity by its
complex conjugate and sum over \(x\in X_m\):
\[
  \Sigma_m
  =
  \frac{1}{F_{m+2}^2}
  \sum_{a,b=0}^{F_{m+2}-1}
  \left(
    \prod_{j=1}^m(1+\zeta^{aF_{j+1}})
  \right)
  \overline{
    \left(
      \prod_{j=1}^m(1+\zeta^{bF_{j+1}})
    \right)
  }
  \sum_{x\in X_m}\zeta^{(b-a)r_x}.
\]
Because \(V_m\) is a bijection onto \(\{0,\dots,F_{m+2}-1\}\), the residues
\(\{r_x:x\in X_m\}\) are exactly the full residue system.  Hence the character
orthogonality relation
\[
  \sum_{x\in X_m}\zeta^{(b-a)r_x}
  =
  \sum_{r=0}^{F_{m+2}-1}\zeta^{(b-a)r}
  =
  \begin{cases}
    F_{m+2}, & a=b,\\
    0, & a\neq b
  \end{cases}
\]
eliminates the cross terms and gives
\[
  \Sigma_m
  =
  \frac{1}{F_{m+2}}
  \sum_{a=0}^{F_{m+2}-1}
  \left|
    \prod_{j=1}^m(1+\zeta^{aF_{j+1}})
  \right|^2,
\]
as claimed.
\end{proof}

\subsection{Affine transfer to Fibonacci subset sums}

\begin{definition}[Shifted subset-sum spectrum]
\label{def:shifted-spectrum}
Define
\[
  \widetilde N_m(\omega):=\sum_{j=1}^m \omega_jF_j
  \qquad (\omega\in\Omega_m),
\]
and
\[
  \widetilde d_m(r):=
  \card{\{\omega\in\Omega_m:\ \widetilde N_m(\omega)=r\}},
  \qquad 0\le r<F_{m+2}.
\]
Because \(\sum_{j=1}^m F_j=F_{m+2}-1\), no reduction modulo \(F_{m+2}\) is
needed in this definition.
\end{definition}

\begin{definition}[The affine residue permutation]
\label{def:affine-map}
Define \(\tau_m\colon \Omega_m\to\Omega_m\) by
\[
  (\tau_m(\omega))_k=
  \begin{cases}
    \omega_{m+1-k}, & k\equiv m \pmod 2,\\[4pt]
    1-\omega_{m+1-k}, & k\not\equiv m \pmod 2,
  \end{cases}
  \qquad 1\le k\le m.
\]
Define
\[
  C_m:=\sum_{\substack{1\le k\le m\\ k\not\equiv m \,(\mathrm{mod}\,2)}}F_{k+1},
  \qquad
  \pi_m(r)\equiv F_{m+1}r+C_m \pmod{F_{m+2}}.
\]
\end{definition}

\begin{lemma}[Closed form for the translation term]
\label{lem:Cm-closed}
For every \(m\ge 1\),
\[
  C_m=F_{m+1}-1.
\]
\end{lemma}

\begin{proof}
If \(m=2q\), then
\[
  C_m=F_2+F_4+\cdots+F_{2q}=F_{2q+1}-1=F_{m+1}-1.
\]
If \(m=2q+1\), then
\[
  C_m=F_3+F_5+\cdots+F_{2q+1}=F_{2q+2}-1=F_{m+1}-1.
\]
\end{proof}

\begin{theorem}[Affine transfer principle]
\label{thm:affine-transfer}
For every \(m\ge 1\), the map \(\pi_m\) is a permutation of
\(\ZZ/F_{m+2}\ZZ\), and
\[
  \widetilde d_m(r)=d_m^\#(\pi_m(r)),
  \qquad
  d_m^\#(s):=d_m(V_m^{-1}(s)).
\]
Equivalently, the multiset of fiber multiplicities of \(\Fold_m\) agrees with
the coefficient multiset of
\[
  \prod_{j=1}^m(1+z^{F_j})
\]
up to the explicit affine permutation \(\pi_m\).
\end{theorem}

\begin{proof}
We use d'Ocagne's identity
\[
  F_uF_{v+1}-F_{u+1}F_v=(-1)^vF_{u-v}.
\]
Taking \(u=m+1\) and \(v=j-1\) gives
\[
  F_{m+1}F_j-F_{m+2}F_{j-1}=(-1)^{j-1}F_{m+2-j},
\]
hence
\begin{equation}\label{eq:docagne-window}
  F_{m+1}F_j\equiv (-1)^{j-1}F_{m+2-j}\pmod{F_{m+2}}.
\end{equation}

Now
\[
  N_m(\tau_m(\omega))
  =
  \sum_{\substack{1\le k\le m\\ k\equiv m \,(\mathrm{mod}\,2)}}
  \omega_{m+1-k}F_{k+1}
  +
  \sum_{\substack{1\le k\le m\\ k\not\equiv m \,(\mathrm{mod}\,2)}}
  (1-\omega_{m+1-k})F_{k+1}.
\]
Separate the constant term:
\[
  N_m(\tau_m(\omega))
  =
  \sum_{k=1}^m(-1)^{m-k}\omega_{m+1-k}F_{k+1}+C_m.
\]
After the change of variables \(j=m+1-k\),
\[
  N_m(\tau_m(\omega))
  =
  \sum_{j=1}^m(-1)^{j-1}\omega_jF_{m+2-j}+C_m.
\]
Applying \eqref{eq:docagne-window} term by term yields
\[
  N_m(\tau_m(\omega))
  \equiv
  F_{m+1}\sum_{j=1}^m \omega_jF_j + C_m
  =
  F_{m+1}\widetilde N_m(\omega)+C_m
  \pmod{F_{m+2}}.
\]
Thus
\begin{equation}\label{eq:tau-congruence}
  N_m(\tau_m(\omega))\equiv \pi_m(\widetilde N_m(\omega))\pmod{F_{m+2}}.
\end{equation}

Since \(\gcd(F_{m+1},F_{m+2})=1\), multiplication by \(F_{m+1}\) is invertible
modulo \(F_{m+2}\), so \(\pi_m\) is a permutation.  Because \(\tau_m\) is also
a bijection of \(\Omega_m\), for every residue \(r\),
\begin{align*}
  \widetilde d_m(r)
  &=
  \card{\{\omega\in\Omega_m:\ \widetilde N_m(\omega)=r\}} \\
  &=
  \card{
    \left\{
      \omega\in\Omega_m:
      N_m(\tau_m(\omega))\equiv \pi_m(r)\pmod{F_{m+2}}
    \right\}
  }.
\end{align*}
By \Cref{thm:residue-fibers}, the last quantity is precisely
\(d_m^\#(\pi_m(r))\).
\end{proof}

\begin{corollary}[Coefficient-spectrum form of the fibers]
\label{cor:coefficient-spectrum}
As multisets,
\[
  \{d_m(x):x\in X_m\}
  =
  \Bigl\{
    [z^r]\prod_{j=1}^m(1+z^{F_j}) : 0\le r<F_{m+2}
  \Bigr\}.
\]
\end{corollary}

\begin{proof}
The coefficient \([z^r]\prod_{j=1}^m(1+z^{F_j})\) is exactly
\(\widetilde d_m(r)\).  Apply \Cref{thm:affine-transfer}.
\end{proof}

\subsection{The partition-difference formula}

\begin{definition}[Fibonacci partition functions]
\label{def:partition-functions}
Let \(R(n)\) denote the number of partitions of \(n\) into distinct Fibonacci
numbers, where the two unit weights \(F_1=F_2=1\) are not distinguished, and
set \(R(n)=0\) for \(n<0\).  Following the indexed model in which the two unit
weights are distinguished, define
\[
  R^\dagger(n):=R(n)+R(n-1)
  \qquad (n\in\ZZ).
\]
\end{definition}

\begin{theorem}[Partition-difference formula]
\label{thm:partition-difference}
For every integer \(n\) with \(0\le n<F_{m+2}\),
\[
  d_m^\#(\pi_m(n))
  =
  R^\dagger(n)-R^\dagger(n-F_{m+1}).
\]
Equivalently,
\[
  d_m^\#(\pi_m(n))
  =
  R(n)+R(n-1)-R(n-F_{m+1})-R(n-F_{m+1}-1).
\]
\end{theorem}

\begin{proof}
By \Cref{thm:affine-transfer},
\[
  d_m^\#(\pi_m(n))=\widetilde d_m(n),
\]
where \(\widetilde d_m(n)\) is the coefficient of \(z^n\) in
\[
  \mathcal P_m(z):=\prod_{j=1}^m(1+z^{F_j}).
\]
Thus
\begin{equation}\label{eq:shifted-coefficient}
  d_m^\#(\pi_m(n))=[z^n]\mathcal P_m(z).
\end{equation}

The ordinary Fibonacci partition function has generating series
\[
  \sum_{n\ge 0}R(n)z^n=(1+z)\prod_{j\ge 3}(1+z^{F_j}),
\]
because the unit weight occurs only once.  Hence
\begin{equation}\label{eq:indexed-generating}
  \sum_{n\ge 0}R^\dagger(n)z^n
  =
  (1+z)\sum_{n\ge 0}R(n)z^n
  =
  \prod_{j\ge 1}(1+z^{F_j}).
\end{equation}
If \(0\le n<F_{m+2}\), then the coefficient of \(z^n\) in the infinite product
\(\prod_{j\ge 1}(1+z^{F_j})\) depends only on the factors with \(j\le m+1\),
since \(F_{m+2}\) and all later Fibonacci numbers exceed \(n\).  Therefore
\[
  R^\dagger(n)
  =
  [z^n](1+z^{F_{m+1}})\mathcal P_m(z)
  =
  [z^n]\mathcal P_m(z)+[z^{n-F_{m+1}}]\mathcal P_m(z).
\]
The second coefficient is \(R^\dagger(n-F_{m+1})\), since
\(n-F_{m+1}<F_{m+1}<F_{m+2}\).  Using \eqref{eq:shifted-coefficient}, we
obtain
\[
  d_m^\#(\pi_m(n))
  =
  R^\dagger(n)-R^\dagger(n-F_{m+1}),
\]
where the second term is interpreted as \(0\) for \(n<F_{m+1}\).  Expanding
\(R^\dagger(t)=R(t)+R(t-1)\) yields the four-term formula.
\end{proof}

\begin{example}[A four-digit window]
\label{ex:m4-partition-difference}
Take \(m=4\).  Then \(F_5=5\), \(F_6=8\), and
\[
  \pi_4(n)\equiv 5n+4 \pmod 8.
\]
Moreover,
\[
  \prod_{j=1}^4(1+z^{F_j})
  =
  (1+z)^2(1+z^2)(1+z^3)
  =
  1+2z+2z^2+3z^3+3z^4+2z^5+2z^6+z^7,
\]
so the difference spectrum is
\[
  \bigl(R^\dagger(n)-R^\dagger(n-5)\bigr)_{n=0}^{7}
  =
  (1,2,2,3,3,2,2,1).
\]
Applying \(\pi_4\) sends \(0,1,\dots,7\) to \(4,1,6,3,0,5,2,7\), and hence
\[
  \bigl(d_4^\#(r)\bigr)_{r=0}^{7}
  =
  (3,2,2,3,1,2,2,1).
\]
Thus the affine permutation simply reindexes the coefficient list of
\(\prod_{j=1}^4(1+z^{F_j})\) into the finite-window fiber spectrum.
\end{example}

\begin{remark}[Single-fiber formulas from classical partition theory]
\label{rem:single-fiber-literature}
Theorem~\ref{thm:partition-difference} identifies each fiber multiplicity with
a Fibonacci-lag difference of the classical partition function.  Hence the
exact Zeckendorf formulas of Chow and Slattery \cite{ChowSlattery2021} apply
directly to \(d_m^\#(\pi_m(n))\).
\end{remark}

\begin{proposition}[Complement symmetry]
\label{prop:complement-symmetry}
Let
\[
  B_m:=\sum_{j=1}^mF_{j+1}=F_{m+3}-2.
\]
Then, as a function on \(\ZZ/F_{m+2}\ZZ\),
\[
  d_m^\#(r)=d_m^\#(B_m-r).
\]
\end{proposition}

\begin{proof}
Let \(\iota(\omega)_j:=1-\omega_j\).  Then \(\iota\) is a bijection of
\(\Omega_m\), and
\[
  N_m(\iota(\omega))
  =
  \sum_{j=1}^m (1-\omega_j)F_{j+1}
  =
  B_m-N_m(\omega).
\]
Hence
\[
  N_m(\omega)\equiv r \pmod{F_{m+2}}
  \iff
  N_m(\iota(\omega))\equiv B_m-r \pmod{F_{m+2}}.
\]
Apply \Cref{thm:residue-fibers}.
\end{proof}
